\begin{abstract}
Skills are increasingly used to extend LLM agents by packaging prompts, code, and configurations into reusable modules. As public registries and marketplaces expand, they form an emerging agentic supply chain, but also introduce a new attack surface for malicious skills. Detecting malicious skills is challenging because relevant evidence is often distributed across heterogeneous artifacts and must be reasoned in context. Existing static, LLM-based, and dynamic approaches each capture only part of this problem, making them insufficient for robust real-world detection. In this paper, we present \toolname, a neuro-symbolic framework for malicious skills detection. \toolname first extracts security-sensitive operations from heterogeneous artifacts through a combination of symbolic parsing and LLM-assisted semantic analysis. It then constructs the skill dependency graph that links artifacts, operations, operands, and value flows across the skill. On top of this graph, \toolname performs neuro-symbolic reasoning to infer malicious patterns or previously unseen suspicious workflows. We evaluate \toolname on a benchmark of 200 real-world skills against 5 state-of-the-art baselines. \toolname achieves 93\% F1, outperforming the baselines by 5\~87 percentage points. We further apply \toolname to analyze 150,108 skills collected from 7 public registries, identifying \textbf{XX} previously unknown malicious skills, all of which were responsibly reported and are currently awaiting confirmation from the platforms and maintainers. These results show that \toolname is both effective and practical for securing the emerging agentic supply chain.
\end{abstract}